\label{chap.4}

The first-generation tester (v1) is a deep RL agent. This chapter formulates the testing
problem as an MDP (Section~\ref{sec:v1-mdp}), describes the agent
(Section~\ref{sec:v1-agent}) and its reward, including the LLM progress term
(Section~\ref{sec:v1-reward}), and reports the three studies that tested it: a
pre-registered transfer experiment (Section~\ref{sec:exp1}), a pre-registered contextual
diagnostic (Section~\ref{sec:exp2}), and checks of the LLM components
(Section~\ref{sec:llm-checks}). Section~\ref{sec:v1-lessons} draws the conclusions that
shaped v2.

\section{Testing as a Markov decision process}
\label{sec:v1-mdp}

An \emph{episode} is a walk of at most 40 actions from the application's entry page. A
\emph{run} (or hunt) is a sequence of episodes on one application until a budget of
$B = 2{,}000$ browser interactions is spent. The state $s_t$ is a summary of the episode so
far and the current page; the actions are those the registry offers on the current page
(Section~\ref{sec:env}), so the action set changes at every step; and the reward is
described in Section~\ref{sec:v1-reward}. The agent never sees the seeded faults or the
oracle's verdicts: learning from them would leak the answer key into the policy.

\paragraph{State.} The state vector has 25 numbers: the episode's history (step, distinct
pages, recent repetitions), the page's composition (counts of links, buttons, fields,
forms, by kind), how much of the current form is filled, and the history of the LLM's
progress judgments. It deliberately carries no page text, heading or URL.

\paragraph{Action features.} Each candidate action $a$ is described by a 95-dimensional
vector $\phi(a)$: its type; a sentence embedding of its label (MiniLM, reduced by PCA); its
role inferred from markup (submit, navigate, edit, remove and so on); flags such as being in
the viewport or inside a form; the category of value to type; and how often it has been
taken in the episode.

\section{The agent: an action-conditioned dueling double DQN}
\label{sec:v1-agent}

A fixed output head, one unit per action slot, cannot work here: slot $k$ means a different
control on every page. The agent instead scores each (state, action) pair with one shared
network, $Q(s, \phi(a))$, following the relevance network of He et al.~\cite{he2016drrn}.
The state is encoded by a 64-unit layer and each action by a 96-unit layer; a 128-unit scorer
combines them into an advantage, and a separate state-value head makes the network
dueling. Because one scorer is shared by every action, the network is
permutation-equivariant over the action set and accepts any number of actions.

Training uses Double DQN targets~\cite{vanhasselt2016double}, 3-step
returns with $\gamma = 0.95$, and prioritised
replay ($\alpha = 0.6$, $\beta_0 = 0.4$, a buffer of 30{,}000
transitions, batches of 64, learning rate $3\times10^{-4}$, target update every 200 steps).

\paragraph{Search-guided exploration.} With uniform $\epsilon$-greedy exploration,
pretraining never reached the order-confirmation pages of the development applications, so
the reward for completing a task never paid. Exploratory actions are therefore taken, with
probability 0.75, from a fixed search priority (the same linear score over the action
features that the \code{search\_S} baseline uses), and otherwise uniformly at random. On the held-out application every learning arm
uses one $\epsilon$ schedule (0.50 decaying to 0.05 over half the budget), so arms do not
differ in how much search they receive.

\paragraph{Online adaptation.} For transfer, the pretrained network is kept and a small
residual head, initialised to zero, is added and trained on the new application. The
pretrained agent therefore starts exactly where pretraining left it and adapts from there.

\paragraph{Bandit variant.} The same network with $\gamma = 0$ (no bootstrapping) is a
contextual bandit. It tests whether multi-step credit assignment adds anything.

\section{Reward}
\label{sec:v1-reward}

The reward has an environment part and an optional LLM part:
\begin{equation}
  r_t = r^{\mathrm{env}}_t + \lambda\, r^{\mathrm{LLM}}_t.
\end{equation}

\paragraph{Environment terms.} A small step cost; a novelty bonus that decays with the
visit count of the abstracted state; a bonus for revealing new controls, capped at one
control per step so that a single-page application does not dominate the scale; a
repetition penalty; and one-time bonuses for the deterministic triggers (console error,
failed request, broken navigation). The total is clipped to $[-6, 6]$.

\paragraph{The LLM progress term.} A local LLM (\code{qwen2.5:7b-instruct}, temperature 0)
reads a text rendering of each newly reached page and rates how far it is through a task, as
a potential $\Phi(s) \in [0,1]$. Judgments are cached by page content, so each distinct page
is judged once. The first design was potential-based shaping,
$\lambda(\gamma\Phi(s') - \Phi(s))$. It preserves the optimal policy, and that is why it
failed: once novelty is exhausted, nothing in the optimal policy reaches outcome pages, and
an agent pretrained on one application never reached its confirmation page in 1{,}500
steps. The term actually used is an \emph{episodic progress record},
\begin{equation}
  r^{\mathrm{LLM}}_t = \max\!\Big(0,\; \Phi(s_{t+1}) - \max_{k \le t} \Phi(s_k)\Big),
\end{equation}
which pays only for a new maximum within the episode, so revisiting or looping earns
nothing and the total per episode is at most 1. A sparse \emph{outcome bonus} of 3 pays the
first time in an episode that a page the model places in the "outcome displayed" stage is
reached, because the model rated any form page highly and completing a task added almost
nothing. Neither term sees the seeded faults.

\section{Experiment 1: does pretraining transfer?}
\label{sec:exp1}

\paragraph{Question.} Does pretraining on five applications make verified-fault discovery on
a held-out sixth faster than learning from scratch? The comparison is against
information-matched non-learning baselines, with and without the LLM reward.

\paragraph{Design.} The design was pre-registered: hypotheses, arms, analysis code and
decision rule were committed before any result existed. It is leave-one-application-out
over the six benchmark applications (Chapter~\ref{chap.6}); four applications were never
used during development (the \emph{pristine} folds). Ten arms ran (Table~\ref{tab:exp1}):
uniform random; two fixed search heuristics (\code{search\_A0}, hand-weighted priority;
\code{search\_S}, the feature-based score that guides exploration); the DQN from scratch and
pretrained, each with and without the LLM reward; the pretrained DQN frozen (zero-shot); and
the bandit from scratch and pretrained. Pretraining used 1{,}500 steps per training
application. Each hunt had $B = 2{,}000$ interactions, with seeds 0--2 and a mechanical
top-up rule. The outcome was $T$, interactions to the first verified fault, right-censored
at $B$. The primary statistic was the per-fold paired difference in restricted mean $T$,
summarised by its median across folds with a fold-level bootstrap interval; the stratified
log-rank test was secondary.

\paragraph{Execution.} There were 291 jobs and 0 failures. The LLM arms ran on a local RTX 4050 GPU
and the other arms on four AWS spot instances (about US\$1.6). The analysis was exactly the
pre-registered one.

\begin{table}[H]
\centering
\small
\caption{Experiment 1: discovery per arm over all folds (runs of 2{,}000 interactions).}
\label{tab:exp1}
\begin{tabular}{@{}lrrrr@{}}
\toprule
Arm & Runs & Found a fault & Rate & Restricted mean $T$ \\
\midrule
\code{dqn\_pretrained\_nollm} & 18 & 7 & 0.39 & 1{,}402 \\
\code{bandit\_pretrained} & 24 & 10 & 0.42 & 1{,}483 \\
\code{dqn\_pretrained} & 27 & 9 & 0.33 & 1{,}520 \\
\code{dqn\_scratch} & 27 & 9 & 0.33 & 1{,}610 \\
\code{dqn\_zeroshot} & 18 & 4 & 0.22 & 1{,}664 \\
\code{bandit\_scratch} & 24 & 5 & 0.21 & 1{,}727 \\
\code{dqn\_scratch\_nollm} & 18 & 5 & 0.28 & 1{,}748 \\
\code{search\_S} & 18 & 3 & 0.17 & 1{,}844 \\
\code{random} & 18 & 2 & 0.11 & 1{,}947 \\
\code{search\_A0} & 18 & 1 & 0.06 & 1{,}953 \\
\bottomrule
\end{tabular}
\end{table}

\begin{table}[H]
\centering
\small
\caption{Experiment 1: pre-registered comparisons (negative $\Delta$ = faster).}
\label{tab:exp1-comp}
\begin{tabular}{@{}lrrr@{}}
\toprule
Comparison & Median $\Delta$ & 95\% interval & Log-rank $p$ \\
\midrule
\code{dqn\_pretrained} $-$ \code{dqn\_scratch} (primary) & 0 & $[-749, +290]$ & 0.527 \\
\code{bandit\_pretrained} $-$ \code{bandit\_scratch} & $-280$ & $[-732, +188]$ & 0.086 \\
\code{dqn\_pretrained} $-$ \code{search\_S} & $-169$ & $[-884, 0]$ & 0.087 \\
\code{dqn\_pretrained} $-$ \code{search\_A0} & $-170$ & $[-1{,}212, 0]$ & 0.006 \\
\code{dqn\_pretrained} $-$ \code{random} & $-151$ & $[-1{,}212, 0]$ & 0.022 \\
LLM reward on $-$ off (pretrained DQN) & $+17$ & $[-122, +378]$ & 0.754 \\
Adaptation on $-$ off & $-84$ & $[-430, 0]$ & 0.534 \\
\bottomrule
\end{tabular}
\end{table}

\paragraph{Results.} Pretraining gave no measurable benefit. The primary comparison has a
median difference of 0 interactions, with an interval reaching far on both sides; on the four
pristine folds it was $-310$ ($p = 0.66$). The LLM reward made no detectable difference.
Learned arms found a fault in more runs than search or random (0.33--0.42 against
0.06--0.17), but their fold-level intervals reach zero and the effect rests largely on one
development fold. Faults 9--16 contextual steps deep, on two applications, were never found
by any arm. The oracle produced no verified finding on any clean build.

A post-hoc description of the hunts showed that about 84\% of each learned arm's actions
were the network's own greedy choice. The null result is not the exploration guide masking
the network: the learned $Q$ drives behaviour, yet pretrained and scratch DQN find faults in
the same share of runs (9 of 27 each). What helped was learned inside the application under
test. This matches the literature's prediction for five training applications per
fold~\cite{zhao2024dinodroid}.

\section{Experiment 2: context or memorisation?}
\label{sec:exp2}

\paragraph{Question.} Does the network rank an action by the state it is in, or has it
memorised actions, words or layouts?

\paragraph{Design.} This experiment was also pre-registered, with separate applications,
runs and report. A seeded generator produced 12 variants of one request workflow. They vary
the vocabulary (training words; synonyms, unseen words and pseudo-words held out with zero
token overlap), the structure (training layouts; tabs, a modal and reordered steps held out)
and add a decoy control. Four agent arms differ only in which inputs the network sees: the
full agent; a label-blind agent; a lexical-only agent; and an action-only agent that never
sees the state. Probe sets of states, with each action's effect labelled by actually
executing it, were fixed before any arm ran. There were 228 jobs (12 pretrainings, 216 hunts)
and 0 failures.

\begin{table}[H]
\centering
\small
\caption{Experiment 2: pre-registered verdicts for the full agent.}
\label{tab:exp2}
\begin{tabularx}{\textwidth}{@{}lXl@{}}
\toprule
Hypothesis & Deciding numbers & Verdict \\
\midrule
H-ctx: state--action contextual learning & within-label AUROC 0.51 $[0.44, 0.59]$; rule $\ge 0.70$ & \textbf{not supported} \\
H-act: action memorisation & decoy probe rate 0.31 $[0.00, 0.91]$ & inconclusive \\
H-lex: lexical memorisation & top-1 ratio on unseen vocabulary 0.88 $[0.35, 1.18]$ & inconclusive \\
H-app: application memorisation & top-1 ratio on new structures 0.63 $[0.44, 0.88]$ & inconclusive \\
\bottomrule
\end{tabularx}
\end{table}

\paragraph{Results.} Contextual learning was not supported. Within a label, the network's
ranking of an action barely depends on whether that action helps in the current state
(AUROC 0.51, chance level), and only 7\% of the variance of $Q$ came from the label $\times$
state interaction: the state moves every action's value together and almost never reorders
them. The agent does rank useful actions above chance overall (top-1 accuracy 0.56, against
0.18 for random); this held up under vocabulary changes but not under every layout change
(tabs 0.35, modal 0.58). On the eight held-out variants, the arms that see the state found
the fault in 19--20 of 24 hunts, against 2 of 24 for search and random, and the label-blind
arm did as well as the full one. The memorisation hypotheses stayed inconclusive with three
seeds: on the decoy, one seed nearly always took it and the other two never did.

\section{Checks of the LLM components}
\label{sec:llm-checks}

\paragraph{Is the progress reward rewarding the right thing?} This post-hoc, descriptive
check asked whether the LLM's progress score separates pages where a fault becomes
observable from pages a random walk visits. On the two development applications, where the
prompt was tuned, it separated them perfectly (AUROC 1.00). On the four pristine
applications it did not (0.73, 0.52, 0.47, 0.38). The score rewards the "happy path", while
many faults need deviations (invalid input, an edit after review), and most of what it paid
duplicated the novelty bonus. This agrees with the warning that small LLM judges make noisy
reward models.

\paragraph{Can the LLM label pages reliably?} For v2, an LLM labeller was to classify page
kinds and control roles once per page. It had to pass gates fixed in advance: page-kind
accuracy $\ge 0.85$ and control-role accuracy $\ge 0.80$.

\begin{table}[H]
\centering
\small
\caption{LLM labeller accuracy against ground truth (gates 0.85 / 0.80).}
\label{tab:labeller}
\begin{tabular}{@{}lrrl@{}}
\toprule
Test & Page kind & Control role & Gate \\
\midrule
Stage 0, generated variants & 0.68 & 0.78 & failed \\
Revised labeller, 69 fresh pages & 0.61 & 0.75 & failed \\
Revised labeller, on its own revision data & 0.93 & 0.83 & (would have passed) \\
\bottomrule
\end{tabular}
\end{table}

Both gates failed. The last row shows why the re-test had to use fresh pages: judged on the
data it was revised from, the labeller would have been adopted.

\section{What v1 taught}
\label{sec:v1-lessons}

Four claims about the deep Q-network are easily conflated. Table~\ref{tab:v1-claims} keeps
them apart.

\begin{table}[H]
\centering
\small
\caption{What the v1 experiments establish about the DQN.}
\label{tab:v1-claims}
\begin{tabularx}{\textwidth}{@{}X X l@{}}
\toprule
Claim & Evidence & Status \\
\midrule
The DQN was implemented and trained & pretraining on five applications per fold; 291 jobs in Experiment~1 with no failures & established \\
Its learned values chose the actions & about 84\% of each learned arm's actions were the network's greedy choice & established \\
It learned which action helps in which state & within-label AUROC 0.51, against a pre-registered bar of 0.70; 7\% of $Q$ variance from label $\times$ state & not supported \\
Pretraining transferred to a new application & median difference 0 interactions, $[-749, +290]$ & not established \\
Learning improved discovery over non-learning arms & learned arms found a fault in more runs (0.33--0.42 against 0.06--0.17), but the intervals reach zero and the effect rests on one fold; held-out H4 inconclusive & suggested only \\
\bottomrule
\end{tabularx}
\end{table}

\begin{table}[H]
\centering
\small
\caption{Findings from v1 and the design consequences for v2.}
\label{tab:lessons}
\begin{tabularx}{\textwidth}{@{}XX@{}}
\toprule
Finding & Consequence for v2 \\
\midrule
Pretraining did not transfer; learning inside the application helped & keep learning, but online, inside each application \\
No state--action contextual learning & make the state explicit: page kind, task stage, form status \\
The LLM reward is misaligned with faults; the labeller failed its gates & remove the LLM from the reward and from the decision path; label pages from markup \\
Deep faults were never reached & explore by remembering states and returning to them (Go-Explore) \\
Faults need deviations from the happy path & deviate deliberately, with test-design operators \\
\bottomrule
\end{tabularx}
\end{table}
